\par\noindent\textbf{(1)} 状态1到2是可逆绝热压缩，\(T_2/T_1=(V_1/V_2)^{\gamma-1}=r^{\gamma-1}\)；3到4是可逆绝热膨胀，\(T_4/T_3=(V_3/V_4)^{\gamma-1}=r^{1-\gamma}\)。两段定容给 \(V_2=V_3,\ V_1=V_4\)。
\par\noindent\textbf{(2)} 定容吸热 \(Q_{\rm in}=nC_v(T_3-T_2)\)、放热幅值 \(Q_{\rm out}=nC_v(T_4-T_1)\)。利用 \(T_4=T_3/r^{\gamma-1}\)、\(T_1=T_2/r^{\gamma-1}\) 得 \(Q_{\rm out}=Q_{\rm in}/r^{\gamma-1}\)，故 \(\eta=1-Q_{\rm out}/Q_{\rm in}=1-r^{1-\gamma}\)。
\par\noindent\textbf{(3)} 对压缩比求导 \(\partial_r\eta=(\gamma-1)r^{-\gamma}>0\)。\(r\downarrow1\) 时 \(\eta\to0\)，\(r\to\infty\) 时形式上 \(\eta\to1\)；但实际气体温度、材料耐热与燃烧控制限制允许压缩比，不能将无穷大极限视作真实发动机效率。
\par\noindent\textbf{(4)} 不可逆压缩、热损失、燃烧不完全及温度相关的 \(C_v\) 使上述可逆模型效率降低或结构改变。特别是两段定容热交换本身在有限温差下不可逆，理想Otto效率也不是两个固定温热库间可逆Carnot效率，不能混淆不同的比较基准。
